% Reviewed translation: PDF pages 227--236 / printed pages 213--222.
% The original scan is authoritative; mathematical tokens were checked against it.
\MathBookSourcePage{227}
\subsection{依赖于网格的范数技术}
\label{subsec:incompressible-sfvp-mesh-dependent}

本小节采用 Babuška, Osborn \& Pitkäranta~\cite{ref:babuska-osborn-pitkaranta-5}
引入的依赖于网格的范数, 在一定程度上改进估计
\eqref{eq:incompressible-sfvp-standard-velocity-error}.
这种范数在这里比 $\|\cdot\|_{\widetilde X}$ 更为合适.
引入它们所依据的思想是在 $\mathcal T_h$ 的每个单元 $\kappa$
上作分部积分. 具体地, 令 $\phi\in H^2(\kappa)$
且 $\theta\in H^1(\kappa)$. Green 公式给出
\[
\int_\kappa\operatorname{curl}\phi\cdot
\operatorname{curl}\theta\,dx
=-\int_\kappa\Delta\phi\,\theta\,dx
+\int_{\partial\kappa}\frac{\partial\phi}{\partial n}\theta\,ds,
\]
其中 $n$ 表示 $\kappa$ 的外单位法向量.

\MathBookSourcePage{228}
为了对 $\mathcal T_h$ 的所有单元求和, 定义集合
\[
\Gamma_h
=\bigcup\{\kappa';\ \kappa'\text{ 是 }\kappa\text{ 的边},
\ \kappa\in\mathcal T_h\},
\]
并定义 $\partial\phi/\partial n$ 跨过 $\kappa$ 的边
$\kappa'$ 的跳跃 $S(\partial\phi/\partial n)$.
更具体地, 若 $\kappa'$ 是 $\kappa_1$ 和 $\kappa_2$ 的公共边,
$n_i$ 表示 $\kappa_i$ 的外法向量, 则置
\[
S(\partial\phi/\partial n)|_{\kappa'}
=\partial\phi/\partial n_1+\partial\phi/\partial n_2;
\]
若 $\kappa'\subset\Gamma$, 则置
\[
S(\partial\phi/\partial n)|_{\kappa'}
=\partial\phi/\partial n.
\]
用这些记号, 有
\begin{equation}\label{eq:incompressible-sfvp-elementwise-green}
\begin{split}
(\operatorname{curl}\phi,\operatorname{curl}\theta)
={}&-\sum_{\kappa\in\mathcal T_h}
\int_\kappa\Delta\phi\,\theta\,dx
+\int_{\Gamma_h}S(\partial\phi/\partial n)\theta\,ds,\\
&\forall\phi\in H^1(\Omega),
\quad \phi|_\kappa\in H^2(\kappa)
\quad\forall\kappa\in\mathcal T_h,
\quad\forall\theta\in H^1(\Omega).
\end{split}
\end{equation}

从现在起, 使用 $L^2$ 和 $H^m$ 空间.
恒等式 \eqref{eq:incompressible-sfvp-elementwise-green}
是另一个与问题 \eqref{eq:incompressible-stream-vorticity-pair}
等价的问题 $(Q)$ 的基础.
事实上, 若 Stokes 问题 \eqref{eq:incompressible-sfvp-stokes}
的解 $\mathbf u$ 和右端项 $\mathbf f$ 满足
\[
\mathbf f\in L^2(\Omega)^2,
\qquad
\omega=\operatorname{curl}\mathbf u\in H^1(\Omega),
\]
则容易证明, 满足
\eqref{eq:incompressible-stream-vorticity-pair} 的 $(\psi,\omega)$
也是下列问题 $(Q)$ 的唯一解:
\[
\left\{
\begin{aligned}
-\nu\sum_{\kappa\in\mathcal T_h}
\int_\kappa\Delta\phi\,\omega\,dx
+\nu\int_{\Gamma_h}S(\partial\phi/\partial n)\omega\,ds
&=(\mathbf f,\operatorname{curl}\phi),\\
(\omega,\mu)
+\sum_{\kappa\in\mathcal T_h}\int_\kappa\Delta\psi\,\mu\,dx
-\int_{\Gamma_h}S(\partial\psi/\partial n)\mu\,ds
&=0,
\end{aligned}
\right.
\]
对所有 $\mu\in H^1(\Omega)$ 和所有 $\phi\in\widetilde\Psi$ 成立, 其中
\[
\widetilde\Psi
=\{\phi\in\Phi;\ \phi|_\kappa\in H^2(\kappa)
\ \forall\kappa\in\mathcal T_h\}.
\]

就逼近而言, 注意 $\Phi_h$ 的函数
(参见 \eqref{eq:incompressible-sfvp-standard-fe-spaces})
整体属于 $\Phi$, 并且在每个 $\kappa$ 上特别属于 $H^2(\kappa)$.
因此恒等式 \eqref{eq:incompressible-sfvp-elementwise-green}
对 $\phi_h\in\Phi_h$ 和 $\theta_h\in\Theta_h$ 成立,
从而问题 \eqref{eq:incompressible-sfvp-practical-problem}
等价于下列问题 $(Q_h)$:

求 $\psi_h\in\Phi_h$ 和 $\omega_h\in\Theta_h$, 使
\begin{equation}\label{eq:incompressible-sfvp-mesh-discrete-problem}
\left\{
\begin{aligned}
-\nu\sum_{\kappa\in\mathcal T_h}
\int_\kappa\Delta\phi_h\,\omega_h\,dx
+\nu\int_{\Gamma_h}S(\partial\phi_h/\partial n)\omega_h\,ds
&=(\mathbf f,\operatorname{curl}\phi_h)
&&\forall\phi_h\in\Phi_h,\\
(\omega_h,\mu_h)
+\sum_{\kappa\in\mathcal T_h}
\int_\kappa\Delta\psi_h\,\mu_h\,dx
-\int_{\Gamma_h}S(\partial\psi_h/\partial n)\mu_h\,ds
&=0
&&\forall\mu_h\in\Theta_h.
\end{aligned}
\right.
\end{equation}

\MathBookSourcePage{229}
观察 \eqref{eq:incompressible-sfvp-elementwise-green}
和 \eqref{eq:incompressible-sfvp-mesh-discrete-problem},
可引入下列双线性形式:
\begin{equation}\label{eq:incompressible-sfvp-mesh-form-a}
a_h(\omega,\mu)=(\omega,\mu)
\qquad\forall\omega,\mu\in L^2(\Omega),
\end{equation}
以及
\begin{equation}\label{eq:incompressible-sfvp-mesh-form-b}
\begin{split}
b_h(\theta,\phi)
={}&\sum_{\kappa\in\mathcal T_h}
\int_\kappa\Delta\phi\,\theta\,dx
-\int_{\Gamma_h}S(\partial\phi/\partial n)\theta\,ds,\\
&\forall\theta\in H^1(\Omega),
\quad\forall\phi\in H^1(\Omega),
\quad\phi|_\kappa\in H^2(\kappa)
\ \forall\kappa\in\mathcal T_h.
\end{split}
\end{equation}
下标 $h$ 用来强调所涉及的某些空间依赖于剖分.
用这些双线性形式, 问题
\eqref{eq:incompressible-sfvp-mesh-discrete-problem}
确实表现为问题 $(Q_h)$:

求 $(\omega_h,\psi_h)\in\Theta_h\times\Phi_h$, 使
\begin{equation}\label{eq:incompressible-sfvp-mesh-abstract-problem}
\left\{
\begin{aligned}
a_h(\omega_h,\theta_h)+b_h(\theta_h,\psi_h)&=0
&&\forall\theta_h\in\Theta_h,\\
b_h(\omega_h,\phi_h)&=-\frac1\nu
(\mathbf f,\operatorname{curl}\phi_h)
&&\forall\phi_h\in\Phi_h.
\end{aligned}
\right.
\end{equation}
显然,
\begin{equation}\label{eq:incompressible-sfvp-mesh-form-b-curl}
b_h(\theta,\phi)
=-(\operatorname{curl}\phi,\operatorname{curl}\theta)
\end{equation}
对所有满足 $\phi|_\kappa\in H^2(\kappa)$ 的
$\phi\in H^1(\Omega)$ 和所有 $\theta\in H^1(\Omega)$ 成立, 并且
\begin{equation}\label{eq:incompressible-sfvp-mesh-form-b-laplacian}
b_h(\theta,\phi)=\int_\Omega\Delta\phi\,\theta\,dx
\end{equation}
对所有满足 $(\partial\phi/\partial n)|_\Gamma=0$ 的
$\phi\in H^2(\Omega)$ 和所有 $\theta\in H^1(\Omega)$ 成立.

还要注意, 若 $\theta$ 仅属于 $L^2(\Omega)$,
且 $\theta|_{\Gamma_h}\in L^2(\Gamma_h)$,
则 $b_h(\cdot,\cdot)$ 仍有定义.
这些考虑促使我们为上述双线性形式配以下依赖于网格的半范数:
\begin{equation}\label{eq:incompressible-sfvp-mesh-norm-zero}
\|\theta\|_{0,h}
=\left(\|\theta\|_{0,\Omega}^2
+h\|\theta\|_{0,\Gamma_h}^2\right)^{1/2}
\qquad\forall\theta\in H^1(\Omega),
\end{equation}
以及
\begin{equation}\label{eq:incompressible-sfvp-mesh-norm-two}
\|\phi\|_{2,h}
=\left(
\sum_{\kappa\in\mathcal T_h}|\phi|_{2,\kappa}^2
+\frac1h\|S(\partial\phi/\partial n)\|_{0,\Gamma_h}^2
\right)^{1/2}
\end{equation}
对所有满足 $\phi|_\kappa\in H^2(\kappa)$ 的
$\phi\in H^1(\Omega)$ 成立.
显然,
\[
|a_h(\omega,\theta)|
\leq\|\omega\|_{0,\Omega}\|\theta\|_{0,\Omega},
\qquad
\forall\omega,\theta\in L^2(\Omega),
\]
以及
\[
|b_h(\theta,\phi)|
\leq\sqrt2\|\theta\|_{0,h}\|\phi\|_{2,h}.
\]

本小节余下几乎全部内容都用于研究这些双线性形式和半范数.
特别地, 将看到这些半范数分别是 $\Theta_h$ 和 $\Phi_h$ 上的范数,
而且 $b_h(\cdot,\cdot)$ 关于它们满足一致 inf--sup 条件.
下一个 Lemma 说明 $\|\cdot\|_{0,h}$ 与 $\|\cdot\|_{0,\Omega}$
是 $\Theta_h$ 上一致等价的两个范数.

\begin{Lemma}\label{lem:incompressible-sfvp-mesh-zero-norm-equivalence}
若 $\mathcal T_h$ 是一族一致正则剖分, 则存在与 $h$ 无关的常数
$C>0$, 使
\begin{equation}\label{eq:incompressible-sfvp-mesh-zero-norm-equivalence}
\|\theta_h\|_{0,h}
\leq C\|\theta_h\|_{0,\Omega}
\qquad\forall\theta_h\in\Theta_h.
\end{equation}
\end{Lemma}

\MathBookSourcePage{230}
证明留给读者作为练习. 提示: 证明
\[
\|\theta_h\|_{0,\Gamma_h}^2
\leq\frac Ch\|\theta_h\|_{0,\Omega}^2
\qquad\forall\theta_h\in\Theta_h.
\]

下一个 Lemma 给出 Corollary~\ref{cor:appendix-global-inverse-gradient}
的一个类似结果.

\begin{Lemma}\label{lem:incompressible-sfvp-mesh-two-norm-inverse}
若 $\mathcal T_h$ 是一族一致正则剖分, 则存在与 $h$ 无关的常数
$C>0$, 使
\begin{equation}\label{eq:incompressible-sfvp-mesh-two-norm-inverse}
\|\theta_h\|_{2,h}
\leq\frac Ch|\theta_h|_{1,\Omega}
\qquad\forall\theta_h\in\Theta_h.
\end{equation}
\end{Lemma}

\begin{Proof}
在每个 $\kappa$ 上对 $\operatorname{grad}\theta_h$
应用 $r=p=2$ 的 \eqref{eq:appendix-local-inverse-inequality}, 得
\[
|\theta_h|_{2,\kappa}
\leq\frac{C_1}{h}|\theta_h|_{1,\kappa}
\qquad\forall\kappa\in\mathcal T_h.
\]
因此只需证明
\begin{equation}\label{eq:incompressible-sfvp-jump-inverse}
\|S(\partial\theta_h/\partial n)\|_{0,\Gamma_h}^2
\leq\frac{C_2}{h}|\theta_h|_{1,\Omega}^2
\qquad\forall\theta_h\in\Theta_h.
\end{equation}
令 $\kappa$ 是 $\mathcal T_h$ 的任一单元.
由 \eqref{eq:appendix-change-forward}
和 \eqref{eq:appendix-gradient-transform}, 得
\[
\int_{\partial\kappa}|\partial\theta_h/\partial n|^2\,ds
\leq\int_{\partial\kappa}\|\operatorname{grad}\theta_h\|^2\,ds
\leq C_3\frac{h_\kappa}{\rho_\kappa^2}
\|\operatorname{grad}\widehat\theta_h\|_{0,\partial\widehat\kappa}^2.
\]
但 $\operatorname{grad}\widehat\theta_h$ 属于有限维空间
$P_{l-1}^2$, 因而
\[
\int_{\partial\widehat\kappa}
|\partial\widehat\theta_h/\partial n|^2\,ds
\leq C_4\frac{\sigma_\kappa}{\rho_\kappa}
|\widehat\theta_h|_{1,\widehat\kappa}^2.
\]
于是附录中的尺度关系给出
\eqref{eq:incompressible-sfvp-jump-inverse}.
\end{Proof}

当然, 不可能把
\eqref{eq:incompressible-sfvp-mesh-zero-norm-equivalence}
推广到 $H^1(\Omega)$, 但容易证明下面的关系.

\begin{Lemma}\label{lem:incompressible-sfvp-mesh-trace-bound}
若 $\mathcal T_h$ 一致正则, 则存在与 $h$ 无关的常数 $C>0$, 使
\begin{equation}\label{eq:incompressible-sfvp-mesh-trace-bound}
\|\theta\|_{0,h}
\leq C\left(\|\theta\|_{0,\Omega}^2
+h^2|\theta|_{1,\Omega}^2\right)^{1/2}
\qquad\forall\theta\in H^1(\Omega).
\end{equation}
\end{Lemma}

\begin{Proof}
由 trace Theorem~\ref{thm:trace-operator},
存在与 $\kappa$ 的几何无关的常数 $C_1$, 使
\[
\|\theta\|_{0,\kappa'}^2
\leq C_1h_\kappa\|\widehat\theta\|_{1,\widehat\kappa}^2
\qquad\forall\theta\in H^1(\kappa).
\]
于是附录中的尺度关系和剖分一致正则性给出
\[
\|\theta\|_{0,\kappa'}^2
\leq C_2\sigma^2
\left[\frac1{th}\|\theta\|_{0,\kappa}^2
+h|\theta|_{1,\kappa}^2\right],
\]
从而证明 \eqref{eq:incompressible-sfvp-mesh-trace-bound}.
\end{Proof}

该 Lemma 有下列有用推论.

\MathBookSourcePage{231}
\begin{Corollary}\label{cor:incompressible-sfvp-mesh-two-norm-coercive}
若 $\mathcal T_h$ 一致正则, 则存在与 $h$ 无关的常数 $C>0$, 使
\begin{equation}\label{eq:incompressible-sfvp-mesh-two-norm-coercive}
\|\phi\|_{2,h}
\geq C\left(
|\phi|_{1,\Omega}^2
+\sum_{\kappa\in\mathcal T_h}|\phi|_{2,\kappa}^2
\right)^{1/2}
\qquad\forall\phi\in\widetilde\Psi.
\end{equation}
\end{Corollary}

\begin{Proof}
由 Lemma~\ref{lem:sec3-boundary-poincare}, 只需证明
\[
\|\phi\|_{2,h}\geq C_1|\phi|_{1,\Omega}.
\]
但由 \eqref{eq:incompressible-sfvp-elementwise-green}
和 \eqref{eq:incompressible-sfvp-mesh-form-b},
\[
|\phi|_{1,\Omega}^2
=-b_h(\phi,\phi)
\leq\sqrt2\|\phi\|_{2,h}\|\phi\|_{0,h},
\]
而 \eqref{eq:incompressible-sfvp-mesh-trace-bound}
结合 Lemma~\ref{lem:sec3-boundary-poincare}
即给出 \eqref{eq:incompressible-sfvp-mesh-two-norm-coercive}.
\end{Proof}

Corollary~\ref{cor:incompressible-sfvp-mesh-two-norm-coercive}
意味着空间 $\widetilde\Psi$ 关于范数 $\|\cdot\|_{2,h}$ 是 Hilbert 空间.
下面两个 Lemma 给出空间 $\Theta_h$ 的逼近性质.

\begin{Lemma}\label{lem:incompressible-sfvp-mesh-zero-approximation}
令 $\mathcal T_h$ 是 $\overline\Omega$ 的正则剖分.
对每个实数 $k\in[1,l+1]$, 存在与 $h$ 无关的常数 $C>0$, 使
\begin{equation}\label{eq:incompressible-sfvp-mesh-zero-approximation}
\inf_{\theta_h\in\Theta_h}
\|v-\theta_h\|_{0,h}
\leq Ch^k|v|_{k,\Omega}
\qquad\forall v\in H^k(\Omega).
\end{equation}
\end{Lemma}

\begin{Proof}
若 $k\geq2$, 取
$\theta_h=I_hv$, 其中插值由
\eqref{eq:appendix-lagrange-interpolant-definition} 定义;
否则取 $\theta_h=R_hv$, 其中局部正则化由附录定义.
于是相应的附录误差估计给出
\[
\|v-\theta_h\|_{0,\Omega}
\leq C_1h^k|v|_{k,\Omega}
\qquad\forall v\in H^k(\Omega).
\]
因此只需估计 $\|v-\theta_h\|_{0,\Gamma_h}$.
与 Lemma~\ref{lem:incompressible-sfvp-mesh-trace-bound} 中一样,
\[
\|v-\theta_h\|_{0,\kappa'}^2
\leq C_1h_\kappa
\|\widehat v-\widehat\theta_h\|_{1,\widehat\kappa}^2.
\]
当 $\theta_h=I_hv$ 时,
\[
\|v-I_hv\|_{0,\kappa'}^2
\leq C_2h_\kappa|\widehat v|_{k,\widehat\kappa}^2,
\]
由尺度关系得到
\[
\|v-I_hv\|_{0,\kappa'}^2
\leq C_3\sigma^2h^{2k-1}|v|_{k,\kappa}^2,
\]
这证明 \eqref{eq:incompressible-sfvp-mesh-zero-approximation}.
当 $\theta_h=R_hv$ 时, 使用
Lemma~\ref{lem:incompressible-sfvp-mesh-trace-bound}
证明中得到的不等式以及 $R_h$ 的局部误差估计
\begin{equation}\label{eq:incompressible-sfvp-local-regularization}
\|v-R_hv\|_{0,\kappa}
+h_\kappa|v-R_hv|_{1,\kappa}
\leq C_5h_\kappa^k|v|_{k,\Delta_\kappa}
\qquad\forall v\in H^k(\Delta_\kappa),
\end{equation}
可得
\[
\|v-R_hv\|_{0,\kappa'}^2
\leq C_6\sigma^2h^{2k-1}|v|_{k,\Delta_\kappa}^2.
\]

\MathBookSourcePage{232}
回顾 $\Delta_\kappa$ 表示 $\mathcal T_h$ 中所有与 $\kappa$
共用顶点或边的单元之并. 此外, 由 $\mathcal T_h$ 的正则性,
任一给定单元在各集合 $\Delta_\kappa$ 中出现的最大次数
由一个与 $h$ 和 $\kappa$ 无关的固定常数 $M$ 控制.
因此对 $\Gamma_h$ 的所有线段 $\kappa'$ 求和,
再次得到 \eqref{eq:incompressible-sfvp-mesh-zero-approximation}.
\end{Proof}

\begin{Lemma}\label{lem:incompressible-sfvp-mesh-two-approximation}
令 $\mathcal T_h$ 是 $\overline\Omega$ 的正则剖分.
对每个实数 $k\in[2,l+1]$, 存在与 $h$ 无关的常数 $C>0$, 使
\begin{equation}\label{eq:incompressible-sfvp-mesh-two-approximation}
\inf_{\theta_h\in\Theta_h}
\|v-\theta_h\|_{2,h}
\leq Ch^{k-2}|v|_{k,\Omega}
\qquad\forall v\in H^k(\Omega).
\end{equation}
\end{Lemma}

\begin{Proof}
由于 $k\geq2$, 可取 $\theta_h=I_hv$.
根据附录中的插值估计, 只需证明
\[
\|S(\partial(v-I_hv)/\partial n)\|_{0,\Gamma_h}^2
\leq C_1h^{2k-3}|v|_{k,\Omega}^2
\qquad\forall v\in H^k(\Omega).
\]
与 Lemma~\ref{lem:incompressible-sfvp-mesh-two-norm-inverse}
中一样, 得
\[
\begin{split}
\|\partial(v-\theta_h)/\partial n\|_{0,\partial\kappa}^2
&\leq C_2\frac{\sigma_\kappa}{\rho_\kappa}
\|\operatorname{grad}(\widehat v-\widehat\theta_h)\|_{0,\partial\widehat\kappa}^2\\
&\leq C_3\frac{\sigma_\kappa}{\rho_\kappa}
\left(|\widehat v-\widehat\theta_h|_{1,\widehat\kappa}^2
+|\widehat v-\widehat\theta_h|_{2,\widehat\kappa}^2\right)\\
&\leq C_4\frac{\sigma_\kappa}{\rho_\kappa}
|\widehat v|_{k,\widehat\kappa}^2.
\end{split}
\]
因此
\[
\|\partial(v-\theta_h)/\partial n\|_{0,\partial\kappa}^2
\leq C_5\sigma^4h^{2k-3}|v|_{k,\kappa}^2,
\]
从而得到 \eqref{eq:incompressible-sfvp-mesh-two-approximation}.
\end{Proof}

应用 Lemma~\ref{lem:incompressible-sfvp-mesh-zero-approximation}
和 Lemma~\ref{lem:incompressible-sfvp-mesh-two-approximation},
最佳估计为
\[
\left.
\begin{aligned}
\inf_{\theta_h\in\Theta_h}\|v-\theta_h\|_{0,h}
&\leq Ch^{l+1}|v|_{l+1,\Omega},\\
\inf_{\theta_h\in\Theta_h}\|v-\theta_h\|_{2,h}
&\leq Ch^{l-1}|v|_{l+1,\Omega},
\end{aligned}
\right\}
\qquad\forall v\in H^{l+1}(\Omega).
\]

下面考察使我们能够应用 Subsection~\ref{subsec:abstract-general-result}
中标准逼近理论的椭圆性和 inf--sup 条件.
由 \eqref{eq:incompressible-sfvp-mesh-zero-norm-equivalence},
双线性形式 $a_h$ 在赋予范数 $\|\cdot\|_{0,h}$ 的 $\Theta_h$
上一致椭圆:
\begin{equation}\label{eq:incompressible-sfvp-mesh-ellipticity}
a_h(\theta_h,\theta_h)
\geq\alpha^*\|\theta_h\|_{0,h}^2
\qquad\forall\theta_h\in\Theta_h,
\quad \alpha^*>0\text{ 与 }h\text{ 无关}.
\end{equation}
更重要的是, 下一个 Lemma 证明 $b_h$
在分别赋予范数 $\|\cdot\|_{0,h}$ 和 $\|\cdot\|_{2,h}$
的 $\Theta_h$ 与 $\Phi_h$ 上满足一致 inf--sup 条件.

\MathBookSourcePage{233}
\begin{Lemma}\label{lem:incompressible-sfvp-mesh-inf-sup}
假设 $\Omega$ 是有界凸多边形, 且剖分 $\mathcal T_h$ 一致正则.
则存在与 $h$ 无关的常数 $\beta^*>0$, 使
\begin{equation}\label{eq:incompressible-sfvp-mesh-inf-sup}
\sup_{\theta_h\in\Theta_h}
\frac{b_h(\theta_h,\phi_h)}{\|\theta_h\|_{0,h}}
\geq\beta^*\|\phi_h\|_{2,h}
\qquad\forall\phi_h\in\Phi_h.
\end{equation}
\end{Lemma}

\begin{Proof}
首先回顾
\[
b_h(\theta_h,\phi_h)
=-(\operatorname{curl}\phi_h,\operatorname{curl}\theta_h)
\qquad\forall\phi_h\in\Phi_h,
\ \forall\theta_h\in\Theta_h.
\]
其次, 对给定 $\phi_h\in\Phi_h$, 由下式定义
$\theta_h\in\Theta_h$:
\begin{equation}\label{eq:incompressible-sfvp-mesh-discrete-neumann-relation}
(\theta_h,\mu_h)
=(\operatorname{curl}\phi_h,\operatorname{curl}\mu_h)
\qquad\forall\mu_h\in\Theta_h.
\end{equation}
注意它与 \eqref{eq:incompressible-sfvp-discrete-four-field}
第二个方程的类比.
由 \eqref{eq:incompressible-sfvp-mesh-ellipticity},
已经找到 $\theta_h\in\Theta_h$, 使
\[
-b_h(\theta_h,\phi_h)
\geq\alpha^*\|\theta_h\|_{0,h}^2.
\]
因此, 为导出 \eqref{eq:incompressible-sfvp-mesh-inf-sup},
将证明
\begin{equation}\label{eq:incompressible-sfvp-mesh-phi-theta-bound}
\|\phi_h\|_{2,h}
\leq C\|\theta_h\|_{0,\Omega}.
\end{equation}

\eqref{eq:incompressible-sfvp-mesh-phi-theta-bound}
的证明与 \eqref{eq:incompressible-sfvp-h1-bound} 的证明非常相似,
因为这里 $\phi_h$ 与 $\theta_h$ 具有相同关系.
不过, 把 \eqref{eq:incompressible-sfvp-mesh-discrete-neumann-relation}
解释为 $\Theta_h$ 中齐次 Neumann 问题的离散化更为方便.
这引导我们引入问题
\[
-\Delta\phi(h)=\theta_h\quad\text{in }\Omega,
\qquad
\partial\phi(h)/\partial n=0\quad\text{on }\Gamma
\]
的解 $\phi(h)$, 它在相差一个加法常数的意义下唯一.
在 \eqref{eq:incompressible-sfvp-mesh-discrete-neumann-relation}
中取 $\mu_h=1\in\Theta_h$, 得
\[
\int_\Omega\theta_h\,dx=0,
\]
这恰是该 Neumann 问题的相容性条件.
此外, $\Omega$ 的凸性蕴含
$\phi(h)\in H^2(\Omega)/\mathbb R$, 且
\[
|\phi(h)|_{2,\Omega}
\leq C_1\|\theta_h\|_{0,\Omega}.
\]
与 Remark~\ref{rem:incompressible-sfvp-elementary-h1-proof} 一样, 可得
\[
|\phi_h-\chi_h|_{1,\Omega}
\leq|\phi(h)-\chi_h|_{1,\Omega}
\qquad\forall\chi_h\in\Theta_h.
\]
由 \eqref{eq:incompressible-sfvp-mesh-two-norm-inverse},
\[
\|\phi_h-\chi_h\|_{2,h}
\leq C_2\frac1h|\phi(h)-\chi_h|_{1,\Omega}.
\]
因此
\[
\|\phi_h\|_{2,h}
\leq\inf_{\chi_h\in\Theta_h}
\left\{
C_2\frac1h|\phi(h)-\chi_h|_{1,\Omega}
+\|\chi_h-\phi(h)\|_{2,h}
+\|\phi(h)\|_{2,h}
\right\}.
\]
取 $\chi_h=I_h\phi(h)$.

\MathBookSourcePage{234}
由 Lemma~\ref{lem:incompressible-sfvp-mesh-two-approximation},
\[
\|\chi_h-\phi(h)\|_{2,h}
\leq C_3|\phi(h)|_{2,\Omega},
\]
而附录插值估计给出
\[
|\chi_h-\phi(h)|_{1,\Omega}
\leq C_4h|\phi(h)|_{2,\Omega}.
\]
此外, 因为 $\phi(h)\in H^2(\Omega)$ 且
$\partial\phi(h)/\partial n=0$ 于 $\Gamma$ 上,
有 $S(\partial\phi(h)/\partial n)=0$, 因而
\[
\|\phi(h)\|_{2,h}=|\phi(h)|_{2,\Omega}.
\]
汇总这些不等式, 得
\eqref{eq:incompressible-sfvp-mesh-phi-theta-bound}.
\end{Proof}

\begin{Remark}\label{rem:incompressible-sfvp-mesh-w1s-inf-sup}
在上述证明中, 由于 $\theta_h$ 与 $\phi_h$
由 \eqref{eq:incompressible-sfvp-mesh-discrete-neumann-relation} 联系,
Lemma~\ref{lem:incompressible-sfvp-hypothesis-h1} 给出
\[
|\phi_h|_{1,s,\Omega}
\leq C_s\|\theta_h\|_{0,\Omega}
\qquad\forall\text{ real }s\geq2.
\]
因此, 在 Lemma~\ref{lem:incompressible-sfvp-mesh-inf-sup}
的假设下, 除 \eqref{eq:incompressible-sfvp-mesh-inf-sup} 外还有
\[
\sup_{\theta_h\in\Theta_h}
\frac{b_h(\theta_h,\phi_h)}{\|\theta_h\|_{0,h}}
\geq\gamma(s)|\phi_h|_{1,s,\Omega}
\qquad\forall\text{ real }s\geq2,
\]
其中常数 $\gamma(s)>0$ 与 $h$ 无关.
\end{Remark}

现在可把问题 \eqref{eq:incompressible-sfvp-mesh-discrete-problem}
置于 Subsection~\ref{subsec:incompressible-abstract-mixed-discretization}
的抽象框架中, 但直接作误差分析实际上更快.
引入与表述 \eqref{eq:incompressible-sfvp-mesh-abstract-problem}
自然相联系的空间
\[
V_h
=\{(\theta_h,\phi_h)\in\Theta_h\times\Phi_h;
\ a_h(\theta_h,\mu_h)+b_h(\mu_h,\phi_h)=0
\ \forall\mu_h\in\Theta_h\}.
\]
应用 Subsection~\ref{subsec:abstract-general-result} 的论证,
容易导出下列逼近结果:
\begin{equation}\label{eq:incompressible-sfvp-mesh-abstract-errors}
\begin{aligned}
\|\omega-\omega_h\|_{0,h}
&\leq(1+\sqrt2C/\alpha^*)
\inf_{(\theta_h,\phi_h)\in V_h}
\|\omega-\theta_h\|_{0,h},\\
\|\omega-\omega_h\|_{0,h}
&\leq(1+\sqrt2C/\alpha^*)
\left\{(1+1/\alpha^*)
\inf_{\theta_h\in\Theta_h}\|\omega-\theta_h\|_{0,h}
+(\sqrt2/\alpha^*)
\inf_{\phi_h\in\Phi_h}\|\psi-\phi_h\|_{2,h}\right\},\\
\|\psi-\psi_h\|_{2,h}
&\leq(1+\sqrt2/\beta^*)
\inf_{\phi_h\in\Phi_h}\|\psi-\phi_h\|_{2,h}
+\frac1{\beta^*}\|\omega-\omega_h\|_{0,\Omega},
\end{aligned}
\end{equation}
其中 $\alpha^*$, $\beta^*$ 和 $C$ 分别是
\eqref{eq:incompressible-sfvp-mesh-ellipticity},
\eqref{eq:incompressible-sfvp-mesh-inf-sup}
和 \eqref{eq:incompressible-sfvp-mesh-phi-theta-bound} 中的常数.
结合 Lemma~\ref{lem:incompressible-sfvp-mesh-zero-approximation}
和 Lemma~\ref{lem:incompressible-sfvp-mesh-two-approximation},
立即得到下一个 Theorem.

\begin{Theorem}\label{thm:incompressible-sfvp-mesh-error}
令 $\Omega$ 是有界凸多边形,
$\mathcal T_h$ 是 $\overline\Omega$ 的一族一致正则剖分.
若 Stokes 问题 \eqref{eq:incompressible-sfvp-stokes}
的解 $\mathbf u=\operatorname{curl}\psi$ 和右端项 $\mathbf f$ 满足
\[
\mathbf f\in L^2(\Omega)^2,
\qquad
\psi\in H^{k+1}(\Omega),
\qquad
\Delta\psi\in H^k(\Omega),
\]
其中某个实数 $k\in[1,l]$, 则问题
\eqref{eq:incompressible-sfvp-mesh-discrete-problem}
的解 $(\omega_h,\psi_h)$ 满足
\begin{equation}\label{eq:incompressible-sfvp-mesh-error}
\|\omega-\omega_h\|_{0,h}
+\|\psi-\psi_h\|_{2,h}
\leq C\left(h^k|\Delta\psi|_{k,\Omega}
+h^{k-1}|\psi|_{k+1,\Omega}\right).
\end{equation}
最佳估计在 $\psi\in H^{l+1}(\Omega)$ 时得到:
\begin{equation}\label{eq:incompressible-sfvp-mesh-best-error}
\|\omega-\omega_h\|_{0,h}
+\|\psi-\psi_h\|_{2,h}
\leq Ch^{l-1}|\psi|_{l+1,\Omega}.
\end{equation}
\end{Theorem}

\MathBookSourcePage{235}
最后, 一个简单的对偶论证可在 $H^1$ 范数中改进 $\psi$ 的误差界.

\begin{Theorem}\label{thm:incompressible-sfvp-mesh-h1-error}
在 Theorem~\ref{thm:incompressible-sfvp-mesh-error} 的假设下,
若流函数 $\psi\in H^{k+1}(\Omega)$,
其中实数 $k\in[2,l]$, 则
\begin{equation}\label{eq:incompressible-sfvp-mesh-h1-error}
|\psi-\psi_h|_{1,\Omega}
\leq Ch^k|\psi|_{k+1,\Omega}.
\end{equation}
当 $\psi\in H^{l+1}(\Omega)$ 时, 最佳误差界为
\begin{equation}\label{eq:incompressible-sfvp-mesh-best-h1-error}
|\psi-\psi_h|_{1,\Omega}
\leq Ch^l|\psi|_{l+1,\Omega},
\qquad l\geq2.
\end{equation}
\end{Theorem}

\begin{Proof}
对 $\mathbf g\in L^2(\Omega)^2$, 考虑辅助 Stokes 问题
\[
\left\{
\begin{aligned}
a_h(\lambda_g,\mu)+b_h(\mu,\phi_g)&=0
&&\forall\mu\in H^1(\Omega),\\
b_h(\lambda_g,\chi)&=-(\mathbf g,\operatorname{curl}\chi)
&&\forall\chi\in\widetilde\Psi.
\end{aligned}
\right.
\]
由于 $\Omega$ 是凸的, 该问题的唯一解
$(\lambda_g,\phi_g)$ 属于 $H^1(\Omega)\times H^3(\Omega)$,
且由 Remark~\ref{rem:stokes-polyhedral-regularity},
\[
\|\phi_g\|_{3,\Omega}+\|\lambda_g\|_{1,\Omega}
\leq C_1\|\mathbf g\|_{0,\Omega}.
\]
由此和问题 \eqref{eq:incompressible-sfvp-mesh-discrete-problem},
容易得到
\[
\begin{split}
(\mathbf g,\operatorname{curl}(\psi-\psi_h))
={}&-b_h(\lambda_g-\lambda_h,\psi-\psi_h)
+a_h(\omega-\omega_h,\lambda_h)\\
={}&b_h(\lambda_h-\lambda_g,\psi-\psi_h)
+a_h(\omega-\omega_h,\lambda_h-\lambda_g)\\
&-b_h(\omega-\omega_h,\phi_g-\phi_h)
\end{split}
\]
对所有 $\lambda_h\in\Theta_h$ 和 $\phi_h\in\Phi_h$ 成立.
于是 \eqref{eq:incompressible-sfvp-mesh-h1-error}
由 Theorem~\ref{thm:incompressible-sfvp-mesh-error},
Lemma~\ref{lem:incompressible-sfvp-mesh-zero-approximation}
和 Lemma~\ref{lem:incompressible-sfvp-mesh-two-approximation} 得到.
\end{Proof}

作为一个有用的推论, 有 $\psi$ 在 $W^{1,s}$ 范数中的下列估计.

\begin{Corollary}\label{cor:incompressible-sfvp-mesh-w1s-error}
假设 Theorem~\ref{thm:incompressible-sfvp-mesh-error} 的条件成立.
若流函数 $\psi\in H^{k+1}(\Omega)$,
其中实数 $k\in[2,l]$, 则对每个实数 $s\geq2$,
存在与 $h$ 无关的常数 $C_s$, 使
\begin{equation}\label{eq:incompressible-sfvp-mesh-w1s-error}
|\psi-\psi_h|_{1,s,\Omega}
\leq C_sh^{k-1+2/s}|\psi|_{k+1,\Omega}.
\end{equation}
\end{Corollary}

\MathBookSourcePage{236}
该估计是 \eqref{eq:incompressible-sfvp-mesh-h1-error},
Lemma~\ref{lem:appendix-global-inverse-same-order}
和附录插值估计的直接推论.

最后三个结果需要作若干评论.
第一个显然的评论是, 本小节的方法给出了比上一小节更尖锐的估计.
事实上, 只要 $l\geq2$ 且流函数 $\psi\in H^{l+1}(\Omega)$,
\eqref{eq:incompressible-sfvp-mesh-best-h1-error}
就给出 $|\psi-\psi_h|_{1,\Omega}$ 的最优误差界.
遗憾的是, 本小节没有改进涡量 $\omega$ 的误差估计.
而且, 当使用分片线性元, 即 $l=1$ 时,
它同样不能证明收敛.
在这种情形下, 精度损失来自范数 $\|\cdot\|_{2,h}$ 中的逼近误差;
显然, 该范数不适合分片线性元.

最后值得指出, 在本小节中几乎不可能去掉 $\Omega$ 的凸性假设,
因为 Lemma~\ref{lem:incompressible-sfvp-mesh-inf-sup}
必然要求 $\phi(h)\in H^2(\Omega)$.

\begin{Remark}\label{rem:incompressible-sfvp-mesh-trace-errors}
由于 $\|\cdot\|_{0,h}$ 和 $\|\cdot\|_{2,h}$ 中含有 trace 项,
\eqref{eq:incompressible-sfvp-mesh-error} 还给出附加估计
\[
\|\omega-\omega_h\|_{0,\Gamma_h}
\leq Ch^{k-3/2}|\psi|_{k+1,\Omega},
\]
以及
\[
\|S(\partial\psi_h/\partial n)\|_{0,\Gamma_h}
\leq Ch^{k-1/2}|\psi|_{k+1,\Omega},
\]
其中 $1\leq k\leq l$, 当然还假设
$\psi\in H^{k+1}(\Omega)$.
最后一个不等式表明, $\partial\psi_h/\partial n$
跨单元边界的跳跃, 以及它在 $\Gamma$ 上的值,
以 $h^{k-1/2}$ 的阶趋于零.
\end{Remark}
